\chapter{Neural Networks for Regression}

\section{From Linear Models to Networks}\index{neural network!regression}
A linear regression predicts $\hat{y} = \mathbf{w}^T\mathbf{x} + b$. A neural network stacks such predictions, with a nonlinear activation between the
layers, so it can fit curved relationships. Think of a curve on a graph. A single straight line cannot follow it, but several bent pieces can. Each \newterm{neuron} computes a weighted sum of its inputs plus a bias, then applies an activation function.

\section{Layers and Activations}\index{activation function}
The input layer holds the features. Hidden layers transform them, and the output layer produces the prediction. Common activations include the
ReLU, $\max(0, z)$, the sigmoid, and the hyperbolic tangent. For regression, the output layer usually uses the identity activation, because the
prediction can be any real number.

\section{Training}
Training a regression network follows the same steps as any gradient-based model.
\begin{itemize}
\item \textbf{Feedforward}: compute the prediction from the inputs.
\item \textbf{Loss}: compare the prediction with the target, usually with mean squared error.
\item \textbf{Backpropagation}: compute the gradient of the loss with respect to each weight.
\item \textbf{Update}: move each weight against its gradient, as in gradient descent.
\end{itemize}

\section{A Worked Example}
Use a single neuron with the identity activation, $\hat{y} = 2x + 1$, on three points $(1, 3)$, $(2, 5)$, and $(3, 8)$. The predictions are
$3, 5, 7$. The errors are $0, 0, 1$, so the mean squared error is
$$\text{MSE} = \frac{0^2 + 0^2 + 1^2}{3} = \frac{1}{3}$$
A neuron with an identity activation is just a linear regression. The power of the network comes from hidden layers with nonlinear activations.

\begin{Exercise}[label=nnr1, title=Compute the error]
A neuron predicts $\hat{y} = 2x + 1$. For the point $(2, 4)$, what is the squared error?
\end{Exercise}
\begin{Answer}[ref=nnr1]
$\hat{y} = 5$, so the squared error is $(5 - 4)^2 = 1$.
\end{Answer}

\begin{Exercise}[label=nnr2, title=Output activation]
Why is the identity activation used in the output layer for regression, but not the sigmoid?
\end{Exercise}
\begin{Answer}[ref=nnr2]
The sigmoid output lies between 0 and 1, which would prevent the network from predicting values outside that range. Regression targets can be any real number, so the identity activation is used.
\end{Answer}
